\documentclass{article}
\usepackage{amsmath,amsfonts}
\usepackage{graphicx}
\usepackage{esint}

\begin{document}

  \title{Derivation of Buoyancy Formula}
  \author{Huang Zixuan 20250065}
  \date{2026.5.22}

  \maketitle

  \par Assume that we have a rigid body $\Sigma$ immersed into the liquid with density $\rho$. Then we partition the boundary of the body $\partial \Sigma$ into infinite surface elements $dS$,
  and also partition the whole body into volume elements $dV$.
  \par For each surface element, the pressure can be expressed as:
  $$p=p_0+\rho gh$$
  by the principle of Passca.
  Accordingly we can express the force:
  $$dF=p\,dS$$
  As the direction of the outer normal vector of surface element and the direction of the force is opposite, we have the vector form:
  $$d\mathbf{F}=-p\cdot d\mathbf{S}$$
  Then the total force can be calculated by integration:
  $$F=\oiint_{\partial \Sigma}dF=\oiint_{\partial \Sigma}-p\cdot d\mathbf{S}$$
  To make the integral match the vector form, we multiply a constant unit vector $\mathbf{c}$ in the integrand:
  $$F=-\oiint_{\partial \Sigma}p\cdot \mathbf{c}\,d\mathbf{S}$$
  Then we consider the vector field $\mathbf{f}=p\cdot \mathbf{c}$, then use Gauss theorem:
  $$F=-\Phi(\mathbf{f},\partial \Sigma)=-\iiint_{\Sigma}\nabla \cdot \mathbf{f}\,dV$$
  As the divergence:
  $$\nabla \cdot \mathbf{f}=\nabla \cdot (p\cdot\mathbf{c})=\nabla p\cdot \mathbf{c}+p(\nabla \cdot \mathbf{c})=\nabla p\cdot \mathbf{c}$$
  As the gradient of $p$ can be computed:
  $$\nabla p=(0,0,\frac{\partial p}{\partial z})=-\rho \mathbf{g}$$
  We have:
  $$F=-\iiint_{\Sigma}-\rho \mathbf{g}\cdot \mathbf{c}\,dV=\rho g\iiint_{\Sigma}dV=\rho gV$$ 
\end{document}